% Source: https://artofproblemsolving.com/community/c3h1599034_mock_sprint_27_semicircles
% Stable identity: aops-topic-1599034-q01
Let $O$ be the centre of the equilateral triangle of centres. Two semicircles of radius $1$ that are externally tangent have centres $2$ apart, so the triangle has side $2$ and circumradius $OC_i=\dfrac{2}{\sqrt3}$.

Each diameter is perpendicular to $OC_i$ at its midpoint $C_i$, so by Pythagoras each endpoint is at distance
\[
\sqrt{\left(\dfrac{2}{\sqrt3}\right)^2+1^2}=\sqrt{\dfrac{7}{3}}
\]
from $O$. All six endpoints are at this distance, so the circle has centre $O$ and radius $R=\sqrt{\dfrac{7}{3}}$.

\begin{asy}
usepackage("amsmath");
size(6cm);
defaultpen(fontsize(10pt));
real s3=sqrt(3);
pair O=(0,0);real R=sqrt(7/3);
fill(circle(O,R),lightgray);
for(int i=0;i<3;++i){
  real ang=90+120*i;pair Ci=(2/s3)*dir(ang);pair u=dir(ang+90);
  pair P1=Ci+u,P2=Ci-u;
  path sc=arc(Ci,1,ang+90,ang+270)--cycle;
  fill(sc,white);draw(arc(Ci,1,ang+90,ang+270),blue+linewidth(1));draw(P1--P2,blue+linewidth(1));
  dot(Ci,blue);dot(P1);dot(P2);
}
draw(circle(O,R),linewidth(1));
pair C0=(2/s3)*dir(90),C1=(2/s3)*dir(210),C2=(2/s3)*dir(330);
draw(C0--C1--C2--cycle,red+dashed);
dot(O);label("$O$",O,S);
pair Pp=C0+dir(180);
draw(O--Pp,dashed);label("$R$",(-1.08,0.6));
label("$1$",C0+0.5*dir(0),N);
label("\mbox{}",(min(currentpicture,user=true)+max(currentpicture,user=true))/2,black); // keeps the text after the figure black when LaTeX typesets the labels inline
\end{asy}

\textbf{The configuration works.} The point where the full circles about $C_i$ and $C_j$ touch is the midpoint of $C_iC_j$, which lies on the inner side of both diameters, so the inward-bulging semicircles do touch there. The semicircles do not overlap one another (their full circles only touch), and they lie inside the big circle: every point of a semicircle is within $1$ of $C_i$, hence within $\dfrac{2}{\sqrt3}+1$ of $O$, but more precisely a point of the half-disc on the inner side of the diameter is at distance at most $\sqrt{\dfrac{7}{3}}$ from $O$ (the farthest points of the half-disc from $O$ are the endpoints of the diameter).

So the shaded area is
\[
\pi\cdot\dfrac{7}{3}-3\cdot\dfrac{\pi}{2}=\dfrac{14\pi-9\pi}{6}=\boxed{\dfrac{5\pi}{6}}.
\]
